\section{The Router Cannot Help You}
\label{sec:ch14-the-router-cannot-help-you}

\index{router!subgraph error propagation|(}%
The obvious next move is to go looking for a switch. The router is the thing
turning four services into one response, so the router ought to be the place to
say what happens when one of them is missing. It has a configuration block with
almost exactly the right name.

\index{router!the keys the block accepts}%
\code{subgraph\_error\_propagation} is real, and
chapter~\ref{ch:enter-the-router}'s trick finds its shape without any
documentation: the router validates \code{config.yaml} against a JSON schema
and refuses to start on a property that schema does not declare, naming the
path it rejected. Feed it an invented key under that block and it names the
block. The keys it does accept cover which errors reach a client and how much
of each one survives the trip: a \code{mode} of \code{wrapped} or
\code{pass-through}, whether the failing service is named, whether locations
and extensions are kept, which extension fields are allowed through, and
whether a subgraph's HTTP status is passed along.

\index{router!what the block does not reach}%
None of it changes anything in this chapter. Run the two requests from
section~\ref{sec:ch14-the-response-that-is-not-there} with the mode left at its
default and again with \code{pass-through}, against the same stopped Ratings
service, and the two responses are byte for byte identical.

\index{router!errors a subgraph returns, not errors about a subgraph}%
The reason is in the block's own name, read carefully. It propagates errors
\emph{from} a subgraph: the ones a running service puts in its own response
body. A service that is not listening never produces a body, so there is
nothing for the block to reshape and the router writes the only sentence it
can. \enquote{Failed to fetch from Subgraph 'ratings'.} is not a formatted
subgraph error. It is the router reporting that there was no subgraph error
because there was no subgraph.

\index{router!what the block does reach}%
The contrast is easy to produce, because chapter~\ref{ch:entities-authorization}
left a field in this graph that a running service refuses. Ask for
\code{Speaker.email} with no token and the Speakers service answers, in its own
body, with an authorization error. Under the default mode the router wraps that:
the client gets one error naming the subgraph, carrying the original inside its
extensions along with the service name and the status. Under
\code{pass-through} the subgraph's own error surfaces directly, with its path
and its \code{AUTH\_NOT\_AUTHENTICATED} code where a client can branch on them.
That is a real and useful difference, and it is a difference about a service
that answered.

\index{blast radius!not a routing concern}%
So the block is worth setting, and it is worth setting for
chapter~\ref{ch:blame-routing}'s reasons rather than this chapter's. It decides
how much a client learns about a failure. It has nothing to say about how much
of the response that failure destroys, because that was decided by the schema,
before any of these services started.

\index{blast radius!the one place it is decided}%
So there is no operational answer to any of this. A router setting will not fix
it, and neither will a retry, a circuit breaker or a health check. The amount
of a response that one service's outage takes with it is a property of the
composed schema, and the only place to change it is the annotations in the
subgraphs that built it.
\index{router!subgraph error propagation|)}%
